\documentclass[10pt,a4paper]{amsart}
\usepackage{amsmath,amssymb,amsthm,graphicx}
% Retained minimal preamble for the Stacks Project chapter more-algebra.tex (More on Algebra).
% The chapter file pinned-original-source.tex carries only \input{preamble} on its first line, so no
% preamble text lives inside the chapter itself.  The Stacks Project uses one single shared preamble for
% every chapter; the snapshot stages it identically as sources/stacks-chow/preamble.tex and
% sources/stacks-etale-cohomology/preamble.tex.  The block below is copied mechanically, line for line,
% from sources/stacks-chow/preamble.tex lines 14-18 (the Xy-pic setup) and 154-217 (the native theorem
% declarations together with the mathematical macros).  Nothing was added, removed or redefined; the
% wrapper class and the amsmath/amssymb/amsthm/graphicx base packages are supplied by the assembler.
% Chapter: more-algebra.tex -- More on Algebra.
% For commutative diagrams we use Xy-pic
\usepackage[all]{xy}

% We use 2cell for 2-commutative diagrams.
\xyoption{2cell}

% Theorem environments.
%
\theoremstyle{plain}
\newtheorem{theorem}[subsection]{Theorem}
\newtheorem{proposition}[subsection]{Proposition}
\newtheorem{lemma}[subsection]{Lemma}

\theoremstyle{definition}
\newtheorem{definition}[subsection]{Definition}
\newtheorem{example}[subsection]{Example}
\newtheorem{exercise}[subsection]{Exercise}
\newtheorem{situation}[subsection]{Situation}

\theoremstyle{remark}
\newtheorem{remark}[subsection]{Remark}
\newtheorem{remarks}[subsection]{Remarks}

\numberwithin{equation}{subsection}

% Macros
%
\def\lim{\mathop{\mathrm{lim}}\nolimits}
\def\colim{\mathop{\mathrm{colim}}\nolimits}
\def\Spec{\mathop{\mathrm{Spec}}}
\def\Hom{\mathop{\mathrm{Hom}}\nolimits}
\def\Ext{\mathop{\mathrm{Ext}}\nolimits}
\def\SheafHom{\mathop{\mathcal{H}\!\mathit{om}}\nolimits}
\def\SheafExt{\mathop{\mathcal{E}\!\mathit{xt}}\nolimits}
\def\Sch{\mathit{Sch}}
\def\Mor{\mathop{\mathrm{Mor}}\nolimits}
\def\Ob{\mathop{\mathrm{Ob}}\nolimits}
\def\Sh{\mathop{\mathit{Sh}}\nolimits}
\def\NL{\mathop{N\!L}\nolimits}
\def\CH{\mathop{\mathrm{CH}}\nolimits}
\def\proetale{{pro\text{-}\acute{e}tale}}
\def\etale{{\acute{e}tale}}
\def\QCoh{\mathit{QCoh}}
\def\Ker{\mathop{\mathrm{Ker}}}
\def\Im{\mathop{\mathrm{Im}}}
\def\Coker{\mathop{\mathrm{Coker}}}
\def\Coim{\mathop{\mathrm{Coim}}}

% Boxtimes
%
\DeclareMathSymbol{\boxtimes}{\mathbin}{AMSa}{"02}

%
% Macros for moduli stacks/spaces
%
\def\QCohstack{\mathcal{QC}\!\mathit{oh}}
\def\Cohstack{\mathcal{C}\!\mathit{oh}}
\def\Spacesstack{\mathcal{S}\!\mathit{paces}}
\def\Quotfunctor{\mathrm{Quot}}
\def\Hilbfunctor{\mathrm{Hilb}}
\def\Curvesstack{\mathcal{C}\!\mathit{urves}}
\def\Polarizedstack{\mathcal{P}\!\mathit{olarized}}
\def\Complexesstack{\mathcal{C}\!\mathit{omplexes}}
% \Pic is the operator that assigns to X its picard group, usage \Pic(X)
% \Picardstack_{X/B} denotes the Picard stack of X over B
% \Picardfunctor_{X/B} denotes the Picard functor of X over B
\def\Pic{\mathop{\mathrm{Pic}}\nolimits}
\def\Picardstack{\mathcal{P}\!\mathit{ic}}
\def\Picardfunctor{\mathrm{Pic}}
\def\Deformationcategory{\mathcal{D}\!\mathit{ef}}
\usepackage{cleveref}
\crefname{originallemma}{original lemma}{original lemmas}
\Crefname{originallemma}{Original lemma}{Original lemmas}
\makeatletter
\newlabel{lemma-pull-root-uniformizer}{{09EV}{source}{}{sourcepointer-1}{}}
\newlabel{lemma-pull-root-uniformizer@cref}{{[originallemma][0][]09EV}{[1][0][]source}{}{}{}}
\makeatother
\crefname{originallemma}{original lemma}{original lemmas}
\Crefname{originallemma}{Original lemma}{Original lemmas}
\makeatletter
\newlabel{lemma-composition-tame}{{0EXU}{source}{}{sourcepointer-2}{}}
\newlabel{lemma-composition-tame@cref}{{[originallemma][0][]0EXU}{[1][0][]source}{}{}{}}
\makeatother
\crefname{originallemma}{original lemma}{original lemmas}
\Crefname{originallemma}{Original lemma}{Original lemmas}
\makeatletter
\newlabel{lemma-subextension-tame}{{0EXV}{source}{}{sourcepointer-3}{}}
\newlabel{lemma-subextension-tame@cref}{{[originallemma][0][]0EXV}{[1][0][]source}{}{}{}}
\makeatother
\crefname{originallemma}{original lemma}{original lemmas}
\Crefname{originallemma}{Original lemma}{Original lemmas}
\makeatletter
\newlabel{lemma-abhyankar}{{0BRM}{source}{}{sourcepointer-4}{}}
\newlabel{lemma-abhyankar@cref}{{[originallemma][0][]0BRM}{[1][0][]source}{}{}{}}
\makeatother
\crefname{originallemma}{original lemma}{original lemmas}
\Crefname{originallemma}{Original lemma}{Original lemmas}
\makeatletter
\newlabel{lemma-extension-dvrs-formally-smooth}{{09E7}{source}{}{sourcepointer-5}{}}
\newlabel{lemma-extension-dvrs-formally-smooth@cref}{{[originallemma][0][]09E7}{[1][0][]source}{}{}{}}
\makeatother
\title{Three equivalent characterizations of tame ramification via adjoining a root of the uniformizer of exponent invertible in the residue field}
\author{The Stacks Project Authors}
\date{stacks-project:more-algebra.tex:local-snapshot-20261012; GNU FDL 1.2 or later}
\begin{document}
\maketitle
\setcounter{subsection}{1}
\noindent\textit{Editorial note.} All delivered mathematical text is the author's own, copied from the pinned Stacks Project chapter file more-algebra.tex; no statement or proof step was rewritten, reordered or removed. The counted statement is the labeled lemma environment at source lines 34184-34199 (canonical tag 0EXW) and the counted proof is the proof environment at source lines 34201-34240, which immediately follows it and is its own proof.\par\medskip
\noindent\textit{Editorial note.} The proof writes the least common multiple of the ramification indices e\_1, ..., e\_r while the maximal ideals are indexed 1, ..., n; both spellings are copied exactly as printed and are not corrected.\par\medskip
\noindent\textit{Editorial note.} The five references are all to earlier results of this chapter: the structure of the root extension, tameness under composition, tameness for subextensions, Abhyankar's lemma and the formally smooth case. All five are navigation stubs with their roles recorded in the reasons.\par\medskip
\noindent\textit{Editorial note.} The chapter file pinned-original-source.tex carries only an input of the shared preamble on its first line, so no preamble text lives inside the chapter. The retained file sources/stacks-more-algebra/retained-preamble.tex is a mechanical, line-for-line copy of the project's single shared preamble as staged in the snapshot (Xy-pic setup plus the native theorem declarations and mathematical macros). Nothing was added, removed or redefined, and the excerpt serves only as preamble, never as a span coordinate.\par\medskip
\noindent\textit{Editorial note.} The source URL is the source-level Stacks Project URL, as required when more than one card is produced from one source; the per-claim tag is carried by claim\_key and by this note. The counted statement and proof cite no bibliography entry, so this card needs no bibliography.\par\medskip
\noindent\textit{Allowed standard prerequisites (editorial).}\begin{itemize}
\item Adjoining an e-th root of a uniformizer produces an extension of discrete valuation rings of ramification index e, so the fraction field of the new ring is the expected one. invoked at proof lines 34202-34204 and again at proof lines 34218-34221. Identifies the ring obtained by adjoining the root as a discrete valuation ring with the correct fraction field and ramification index.
\item Tameness is preserved by composition of extensions of discrete valuation rings and inherited by subextensions. invoked at proof lines 34205-34207. Deduces tameness of L over K from tameness of the larger extension L' and from the fact that L is a subextension.
\item Abhyankar's lemma: after adjoining a root of the uniformizer of exponent equal to the ramification index, the extension becomes formally smooth in the adic topology, hence unramified. invoked at proof lines 34231-34237. Shows that each factor of the base change becomes unramified over the new base, which is condition (2).
\end{itemize}\medskip
\ifdefined\hypertarget\hypertarget{sourcepointer-1}{}\fi
\noindent\textit{Source reference.} Original lemma 09EV: sources/stacks-more-algebra/pinned-original-source.tex, source lines 33954--33954. The reference is used twice, at proof lines 34202-34203 and 34220-34221, to know that adjoining the root of the uniformizer gives a discrete valuation ring with the expected fraction field and ramification index; the argument built on it is inside the counted span.\par\medskip
\ifdefined\hypertarget\hypertarget{sourcepointer-2}{}\fi
\noindent\textit{Source reference.} Original lemma 0EXU: sources/stacks-more-algebra/pinned-original-source.tex, source lines 34138--34138. The reference is applied at proof lines 34205-34206 to see that the larger extension is tamely ramified over A; the deduction for the subextension is the next step inside the counted span.\par\medskip
\ifdefined\hypertarget\hypertarget{sourcepointer-3}{}\fi
\noindent\textit{Source reference.} Original lemma 0EXV: sources/stacks-more-algebra/pinned-original-source.tex, source lines 34160--34160. The reference is applied at proof line 34207 to pass from tameness of the larger extension to tameness of L; it is one step of the implication (2) implies (1), the rest of which is inside the counted span.\par\medskip
\ifdefined\hypertarget\hypertarget{sourcepointer-4}{}\fi
\noindent\textit{Source reference.} Original lemma 0BRM: sources/stacks-more-algebra/pinned-original-source.tex, source lines 34037--34037. The reference is applied at proof lines 34231-34234 to obtain formal smoothness of each factor over the extended base; the decomposition and the choice of e\_0 that make that application possible are inside the counted span.\par\medskip
\ifdefined\hypertarget\hypertarget{sourcepointer-5}{}\fi
\noindent\textit{Source reference.} Original lemma 09E7: sources/stacks-more-algebra/pinned-original-source.tex, source lines 32989--32989. The reference converts formal smoothness into ramification index one and separable residue extension at proof lines 34235-34237, which is the last step towards condition (2); it is a navigation pointer and the construction that reaches it is inside the counted span.\par\medskip
\begin{lemma}
\label{lemma-characterize-tame}
Let $A$ be a discrete valuation ring with fraction field $K$.
Let $\pi \in A$ be a uniformizer.
Let $L/K$ be a finite separable extension.
The following are equivalent
\begin{enumerate}
\item $L$ is tamely ramified with respect to $A$,
\item there exists an $e \geq 1$ invertible in $\kappa_A$
and an extension $L'/K' = K[\pi^{1/e}]$ unramified with respect to
$A' = A[\pi^{1/e}]$ such that $L$ is contained in $L'$, and
\item there exists an $e_0 \geq 1$ invertible in $\kappa_A$
such that for every $d \geq 1$ invertible in $\kappa_A$
(2) holds with $e = de_0$.
\end{enumerate}
\end{lemma}

\begin{proof}
Observe that $A'$ is a discrete valuation ring with fraction
field $K'$, see Lemma \ref{lemma-pull-root-uniformizer}.
Of course the ramification index of $A'$ over $A$ is $e$.
Thus if (2) holds, then $L'$ is tamely ramified with respect to $A$
by Lemma \ref{lemma-composition-tame}. Hence $L$ is tamely
ramified with respect to $A$ by Lemma \ref{lemma-subextension-tame}.

\medskip\noindent
The implication (3) $\Rightarrow$ (2) is immediate.

\medskip\noindent
Assume that (1) holds. Let $B$ be the integral closure of $A$ in $L$
and let $\mathfrak m_1, \ldots, \mathfrak m_n$ be its maximal ideals.
Denote $e_i$ the ramification index of $A \to B_{\mathfrak m_i}$.
Let $e_0$ be the least common multiple of $e_1, \ldots, e_r$.
This is invertible in $\kappa_A$ by our assumption (1).
Let $e = de_0$ as in (3). Set $A' = A[\pi^{1/e}]$.
Then $A \to A'$ is an extension of discrete valuation rings
with fraction field $K' = K[\pi^{1/e}]$, see
Lemma \ref{lemma-pull-root-uniformizer}.
Choose a product decomposition
$$
L \otimes_K K' = \prod L'_j
$$
where $L'_j$ are fields. Let $B'_j$ be the integral closure
of $A$ in $L'_j$. Let $\mathfrak m_{ijk}$ be the maximal
ideals of $B'_j$ lying over $\mathfrak m_i$.
Observe that $(B'_j)_{\mathfrak m_i}$ is the integral
closure of $B_{\mathfrak m_i}$ in $L'_j$.
By Abhyankar's lemma (Lemma \ref{lemma-abhyankar})
applied to $A \subset B_{\mathfrak m_i}$ and the extension $K'/K$
we see that $A' \to (B'_j)_{\mathfrak m_{ijk}}$
is formally smooth in the $\mathfrak m_{ijk}$-adic topology.
This implies that the ramification index is $1$ and
that the residue field extension is separable
(Lemma \ref{lemma-extension-dvrs-formally-smooth}).
In this way we see that $L'_j$ is unramified with respect to $A'$.
This finishes the proof: we take $L' = L'_j$ for some $j$.
\end{proof}

\end{document}
